\section{Discussion}
\label{sec:discussion}

We have shown that the ladder structure of $\cN=2$ SYK reduces to a single Dirac response on the AdS$_2$ strip, with two scalars as its components in the bosonic reduction, and that a constant double-trace superpotential on the cylinder breaks the bulk superalgebra from $su(1,1|1)$ to $u(1|1)$. The spectrum retains the shape of the full multiplets at all $g$, as in Thong's case where the dimensions and OPE data survive the loss of $E_\pm$. We list what we consider the main open points.

\paragraph{The super-centre integral.} \Cref{subsec:supercentre} maps each superblock to one exponential, and the four-point function to the Dirac propagator, so one Dirac field is the $\cN=2$ replacement of Thong's single scalar at the level of four-point data. Two points remain open there. The continued superfamilies are defined through their components, and a characterization in superspace, for instance through the eigenfunctions~\cite[(4.40)]{Bulycheva:2018qcp} on $z>1$, would be more natural. And the finite part at the collisions is a prescription.

\paragraph{Term-by-term reduction.} We have not proved the analogue of~\cite[App.~C.2]{Thong:2026} for the $\cN=2$ families, ordinary or super, including the step-function families, or that the mixed $\psi$--$b$ continuation allows the required exchanges of integrals. Working with $G(\ket{12})$ directly may avoid the separate infrared and ultraviolet divergences of the fermion and boson rails.

\paragraph{A modified boundary theory.} \Cref{prop:breaking} concerns the given boundary domain. A massive analogue of the Type II condition of~\cite{Correa:2019rdk} might preserve more, at the cost of a derivative term in the scalar condition. It isn't clear that such a theory would reproduce the ladder data.

\paragraph{The physical $H=1$ coefficient.} Connecting~\eqref{eq:finite-part} to the super-Schwarzian~\cite{Fu:2016vas,Forste:2017apw} is a well-defined next computation, which the literature we have read leaves undone~\cite{Peng:2017spg,Bulycheva:2018qcp}.

\paragraph{Antipodal coupling.} The boundary vertex couples the two boundaries of global AdS$_2$ at equal time (\cref{fig:strip}). That has the form of the double-trace couplings that make wormholes traversable~\cite{Gao:2016bin,Maldacena:2018lmt}. We mention this only as a structural similarity. Here the bulk is rigid and the fields are auxiliary, and the two boundaries are the two signs of a relative separation, not two copies of a system.

\paragraph{Stepping stones.} Complex SYK at zero charge is a useful test. We find its symmetric-channel kernel~\cite{Bulycheva:2017uqj,Gu:2019jub} satisfies
\begin{equation}
k^S(h)\,\cR_{\nu,0}(h-\half)=4^{-\nu},\qquad\nu=1-\frac2q,
\end{equation}
to be compared with $k^A\cR_{\nu,1}=(q-1)4^{-\nu}$~\cite{Thong:2026}. So it is one scalar with a parity-dependent Robin coefficient, with a mode at $h=1$ for the $U(1)$ charge and $a+J=2/q-\half<0$ for $q>4$. It is the simplest system where the even-sector issues of \cref{sec:stability} arise without mixing. $\cN=1$ tests something else, because it has an odd relative coordinate.

\paragraph{Beyond this paper.} Nonzero charge~\cite{Gu:2019jub,Heydeman:2022lse} would need a $\theta$-dependent mixing. Higher-point data remain out of reach. Thong's caveat stands: this construction concerns four-point data, not an equivalence of theories.
